\begin{PSbox}[unbreakable]{obj}{Learning Objectives}

After completing this module, students will be able to:

\begin{enumerate}
\def\labelenumi{\arabic{enumi}.}
\tightlist
\item
  Define random experiments and identify sample spaces for engineering scenarios
\item
  Apply set operations to define complex events
\item
  Use probability axioms to derive probability rules
\item
  Compute conditional probabilities
\item
  Determine statistical independence
\item
  Apply Bayes’ theorem to engineering inference problems
\item
  Use the total probability theorem for multi-path analysis
\item
  Simulate random experiments in MATLAB
\end{enumerate}

\end{PSbox}

\PSsection{1.1}{Introduction: Why Probability in Engineering?}

\PSsubsection{The Engineering Reality}

No engineering system operates in a perfectly deterministic world. Consider:

\begin{itemize}
\tightlist
\item
  A digital communication system transmits bits, but \textbf{noise} corrupts some bits randomly
\item
  A manufacturing line produces resistors, but their values \textbf{vary} around the nominal
\item
  A sensor measures temperature, but each reading has \textbf{random error}
\item
  A wireless signal travels to a receiver, but \textbf{fading} causes random attenuation
\end{itemize}

\textbf{Probability theory} provides the mathematical framework to model, analyze, and design systems that operate under uncertainty.

\PSsubsection{Engineering Question}

\begin{PSquote}

``If I transmit 1000 bits over a noisy channel, how many errors should I expect? Can I guarantee fewer than 5 errors?''

\end{PSquote}

To answer this, we need:

\begin{enumerate}
\def\labelenumi{\arabic{enumi}.}
\tightlist
\item
  A model of the random phenomenon (bit errors)
\item
  Mathematical tools to compute probabilities
\item
  Methods to make engineering decisions based on those probabilities
\end{enumerate}

\PSsection{1.2}{Random Experiments}

\begin{PSbox}[unbreakable]{def}{Definition}

A \textbf{random experiment} is a procedure that:

\begin{enumerate}
\def\labelenumi{\arabic{enumi}.}
\tightlist
\item
  Can be repeated under identical conditions
\item
  Has a well-defined set of possible outcomes
\item
  Has an outcome that cannot be predicted with certainty before the experiment is performed
\end{enumerate}

\end{PSbox}

\PSsubsection{Engineering Examples}

{\def\LTcaptype{none} % do not increment counter
\begin{longtable}[c]{@{}
  >{\centering\arraybackslash}p{(0.965\linewidth - 4\tabcolsep) * \real{0.2821}}
  >{\centering\arraybackslash}p{(0.965\linewidth - 4\tabcolsep) * \real{0.2308}}
  >{\centering\arraybackslash}p{(0.965\linewidth - 4\tabcolsep) * \real{0.4872}}@{}}
\toprule\noalign{}
\rowcolor{PSnavy}\PSth{Experiment} & \PSth{Context} & \PSth{Uncertainty Source} \\
\midrule\noalign{}
\endhead
\bottomrule\noalign{}
\endlastfoot
Transmit one bit over a noisy channel & Communications & Thermal noise \\
Measure a resistor’s value & Electronics & Manufacturing variation \\
Count packets arriving in 1 second & Networking & Traffic randomness \\
Measure received signal power & Wireless & Fading, interference \\
Test a component until failure & Reliability & Material defects, wear \\
Sample ADC output voltage & Signal processing & Quantization + noise \\
\end{longtable}
}

\PSsubsection{Key Distinction}

{\def\LTcaptype{none} % do not increment counter
\begin{longtable}[c]{@{}
  >{\centering\arraybackslash}p{(0.965\linewidth - 2\tabcolsep) * \real{0.6364}}
  >{\centering\arraybackslash}p{(0.965\linewidth - 2\tabcolsep) * \real{0.3636}}@{}}
\toprule\noalign{}
\rowcolor{PSnavy}\PSth{Deterministic} & \PSth{Random} \\
\midrule\noalign{}
\endhead
\bottomrule\noalign{}
\endlastfoot
Output of a logic gate with known inputs & Output of a logic gate with noisy inputs \\
Resistance calculated from color code & Actual measured resistance \\
Free-space path loss at known distance & Actual received power with fading \\
\end{longtable}
}

\PSsection{1.3}{Sample Spaces}

\begin{PSbox}[unbreakable]{def}{Definition}

The \textbf{sample space} $S$ (or $\Omega$) is the set of all possible outcomes of a random experiment.

\end{PSbox}

\PSsubsection{Types of Sample Spaces}

\textbf{Discrete (Countable):}

\begin{itemize}
\tightlist
\item
  Transmit one bit: $S = \{0,\, 1\}$ (received correctly or in error)
\item
  Number of errors in a packet of $n$ bits: $S = \{0,\, 1,\, 2,\, \ldots,\, n\}$
\item
  Number of arrivals at a server: $S = \{0,\, 1,\, 2,\, 3,\, \ldots \}$
\end{itemize}

\textbf{Continuous (Uncountable):}

\begin{itemize}
\tightlist
\item
  Measured voltage: $S = \mathbb{R}$ (or a practical range like $[0,\, 5] V$)
\item
  Time until component failure: $S = [0,\, \infty)$
\item
  Phase of received signal: $S = [0,\, 2\pi)$
\end{itemize}

\begin{PSbox}[unbreakable]{ex}{Engineering Example: Communication System}

A transmitter sends a bit (0 or 1) over a noisy channel. The receiver makes a decision.

\PSfigure{m01-sample-space}

\end{PSbox}

\PSsection{1.4}{Events}

\begin{PSbox}[unbreakable]{def}{Definition}

An \textbf{event} is a subset of the sample space. An event \textbf{occurs} if the outcome of the experiment belongs to that subset.

\end{PSbox}

\PSsubsection{Types of Events}

\begin{itemize}
\tightlist
\item
  \textbf{Simple event}: Contains exactly one outcome. $E = \{(0,\,1)\}$
\item
  \textbf{Compound event}: Contains multiple outcomes. $E = \{(0,\,1),\, (1,\,0)\}$ (any error)
\item
  \textbf{Certain event}: The entire sample space $S$ (always occurs)
\item
  \textbf{Impossible event}: The empty set $\emptyset$ (never occurs)
\end{itemize}

\PSsubsection{Set Operations on Events}

{\def\LTcaptype{none} % do not increment counter
\begin{longtable}[c]{@{}
  >{\centering\arraybackslash}p{(0.965\linewidth - 6\tabcolsep) * \real{0.1930}}
  >{\centering\arraybackslash}p{(0.965\linewidth - 6\tabcolsep) * \real{0.1754}}
  >{\centering\arraybackslash}p{(0.965\linewidth - 6\tabcolsep) * \real{0.1579}}
  >{\centering\arraybackslash}p{(0.965\linewidth - 6\tabcolsep) * \real{0.4737}}@{}}
\toprule\noalign{}
\rowcolor{PSnavy}\PSth{Operation} & \PSth{Notation} & \PSth{Meaning} & \PSth{Engineering Interpretation} \\
\midrule\noalign{}
\endhead
\bottomrule\noalign{}
\endlastfoot
Union & $A \cup B$ & A or $B$ (or both) occurs & Either failure mode A or $B$ triggers alarm \\
Intersection & $A \cap B$ & Both A and $B$ occur & Both sensors detect the signal \\
Complement & $A^{c}$ or Ā & A does not occur & Component does NOT fail \\
Difference & $A - B = A \cap B^{c}$ & A occurs but not $B$ & Sensor A triggers but sensor $B$ doesn’t \\
Mutually exclusive & $A \cap B = \emptyset$ & A and $B$ cannot both occur & Component can’t be both working and failed \\
\end{longtable}
}

\begin{PSbox}[unbreakable]{ex}{Engineering Example: Manufacturing Defects}

A circuit board inspection checks for three types of defects:

\begin{itemize}
\tightlist
\item
  A = solder bridge defect
\item
  $B$ = missing component
\item
  $C$ = wrong component value
\end{itemize}

Events:

\begin{itemize}
\tightlist
\item
  $A \cup B \cup C$ = ``board has at least one defect''
\item
  $A \cap B$ = ``board has both a solder bridge AND a missing component''
\item
  $(A \cup B \cup C)^{c}$ = ``board is defect-free''
\end{itemize}

\end{PSbox}

\PSsection{1.5}{Probability Axioms}

\PSsubsection{Kolmogorov’s Axioms}

For a sample space $S$ with events defined on it, a probability function $P$ satisfies:

\textbf{Axiom 1 (Non-negativity):}\PSdisplay{P(A) \geq 0 \quad \text{for any event } A}

\textbf{Axiom 2 (Normalization):}\PSdisplay{P(S) = 1}

\textbf{Axiom 3 (Countable Additivity):}\PSbr{}For mutually exclusive events $A_{1},\, A_{2},\, A_{3}$, …:\PSdisplay{P(A_1 \cup A_2 \cup A_3 \cup \cdots) = P(A_1) + P(A_2) + P(A_3) + \cdots}

\PSsubsection{Properties Derived from Axioms}

\textbf{Property 1: Complement Rule}\PSdisplay{P(A^c) = 1 - P(A)}

\emph{Proof:} $S = A \cup A^{c}$ and $A \cap A^{c} = \emptyset$, so $P(S) = P(A) + P(A^{c}) = 1$.

\emph{Engineering meaning:} $P(\text{component works}) = 1 - P(\text{component fails})$

\textbf{Property 2: Impossible Event}\PSdisplay{P(\emptyset) = 0}

\textbf{Property 3: Monotonicity}\PSbr{}If $A \subseteq B$, then $P(A) \le P(B)$

\textbf{Property 4: Bound}\PSdisplay{0 \leq P(A) \leq 1}

\textbf{Property 5: Inclusion-Exclusion (Two Events)}\PSdisplay{P(A \cup B) = P(A) + P(B) - P(A \cap B)}

\emph{Engineering meaning:} When computing the probability of ``failure $A$ OR failure $B$,'' we must subtract the overlap where both fail simultaneously.

\PSsection{1.6}{Probability Rules}

\PSsubsection{Addition Rule}

For any two events:\PSdisplay{P(A \cup B) = P(A) + P(B) - P(A \cap B)}

For mutually exclusive events $(A \cap B = \emptyset)$:\PSdisplay{P(A \cup B) = P(A) + P(B)}

\PSsubsection{Complement Rule}

\PSdisplay{P(A) = 1 - P(A^c)}

This is often the easiest way to compute probabilities:

\begin{PSquote}

$P(\text{at least one error in 100}\,\mathrm{bits}) = 1 - P(\text{zero errors in 100}\,\mathrm{bits})$

\end{PSquote}

\PSsubsection{Inclusion-Exclusion (Three Events)}

\PSdisplay{P(A \cup B \cup C) = P(A) + P(B) + P(C) - P(A \cap B) - P(A \cap C) - P(B \cap C) + P(A \cap B \cap C)}

\begin{PSbox}[unbreakable]{ex}{Engineering Example: Network Redundancy}

A data center has three independent network links. Each link fails with probability 0.05.

\begin{itemize}
\tightlist
\item
  A = link 1 fails, $B$ = link 2 fails, $C$ = link 3 fails
\item
  $P(\text{at least one fails}) = P(A \cup B \cup C)$
\item
  Using inclusion-exclusion:
\end{itemize}

$P(A \cup B \cup C) = 3(0.05) - 3(0.05)^{2} + (0.05)^{3} = 0.15 - 0.0075 + 0.000125 = 0.142625$

Alternatively: $P(\text{at least one fails}) = 1 - P(\text{none fail}) = 1 - (0.95)^{3} = 0.142625$ \PSyes{}

\end{PSbox}

\PSsection{1.7}{Conditional Probability}

\begin{PSbox}[unbreakable]{def}{Definition}

The \textbf{conditional probability} of event A given that event $B$ has occurred:

\PSdisplay{P(A|B) = \frac{P(A \cap B)}{P(B)}, \quad P(B) > 0}

\end{PSbox}

\PSsubsection{Interpretation}

Conditional probability \textbf{updates} our knowledge. Once we know $B$ occurred, the effective sample space shrinks from $S$ to $B$, and we ask: ``Within $B$, how likely is $A$?''

\PSsubsection{Key Distinction}

{\def\LTcaptype{none} % do not increment counter
\begin{longtable}[c]{@{}
  >{\centering\arraybackslash}p{(0.965\linewidth - 4\tabcolsep) * \real{0.3214}}
  >{\centering\arraybackslash}p{(0.965\linewidth - 4\tabcolsep) * \real{0.3571}}
  >{\centering\arraybackslash}p{(0.965\linewidth - 4\tabcolsep) * \real{0.3214}}@{}}
\toprule\noalign{}
\rowcolor{PSnavy}\PSth{Concept} & \PSth{Notation} & \PSth{Meaning} \\
\midrule\noalign{}
\endhead
\bottomrule\noalign{}
\endlastfoot
Probability of $A$ & $P(A)$ & Likelihood of A with no additional information \\
Conditional probability & $P(A \mid B)$ & Likelihood of $A$ \textbf{given} that $B$ has occurred \\
Joint probability & $P(A \cap B)$ & Likelihood that \textbf{both} A and $B$ occur \\
\end{longtable}
}

These are related: $P(A \cap B) = P(A \mid B) \cdot P(B) = P(B \mid A) \cdot P(A)$

\begin{PSbox}[unbreakable]{ex}{Engineering Example: Detection Theory}

A radar system detects aircraft:

\begin{itemize}
\tightlist
\item
  $H_{1}$ = aircraft is present (target)
\item
  $H_{0}$ = aircraft is absent (no target)
\item
  $D$ = detector says ``target present''
\end{itemize}

Important conditional probabilities:

\begin{itemize}
\tightlist
\item
  $P(D \mid H_{1})$ = \textbf{probability of detection} (sensitivity)
\item
  $P(D \mid H_{0})$ = \textbf{probability of false alarm}
\item
  $P(D^{c} \mid H_{1})$ = \textbf{probability of miss} $= 1 - P(D \mid H_{1})$
\end{itemize}

Suppose: $P(H_{1}) = 0.01,\, P(D \mid H_{1}) = 0.95,\, P(D \mid H_{0}) = 0.02$

What is $P(\text{target actually present}\ \mid \text{detector says present})$?

This requires \textbf{Bayes’ theorem} (Section 1.9).

\end{PSbox}

\begin{PSbox}[unbreakable]{ex}{Engineering Example: Communication Channel}

A binary symmetric channel has crossover probability $p = 0.01$:

\begin{itemize}
\tightlist
\item
  $P(\text{receive 1}\ \mid \text{sent 0}) = p = 0.01$
\item
  $P(\text{receive 0}\ \mid \text{sent 1}) = p = 0.01$
\item
  $P(\text{receive 0}\ \mid \text{sent 0}) = 1 - p = 0.99$
\item
  $P(\text{receive 1}\ \mid \text{sent 1}) = 1 - p = 0.99$
\end{itemize}

If bits are equally likely to be 0 or 1:

\begin{itemize}
\tightlist
\item
  $P(\text{sent 0}\ \cap \text{receive 1}) = P(\text{receive 1}\ \mid \text{sent 0}) \cdot P(\text{sent 0}) = 0.01 \times 0.5 = 0.005$
\item
  $P(\text{error}) = P(\text{sent 0}\ \cap \text{receive 1}) + P(\text{sent 1}\ \cap \text{receive 0}) = 0.005 + 0.005 = 0.01$
\end{itemize}

\end{PSbox}

\PSsubsection{Properties of Conditional Probability}

Conditional probability is itself a valid probability measure:

\begin{enumerate}
\def\labelenumi{\arabic{enumi}.}
\tightlist
\item
  $P(A \mid B) \ge 0$
\item
  $P(S \mid B) = 1$
\item
  $P(A_{1} \cup A_{2} \mid B) = P(A_{1} \mid B) + P(A_{2} \mid B)$ if $A_{1} \cap A_{2} = \emptyset$
\end{enumerate}

\PSsubsection{Multiplication Rule (Chain Rule)}

\PSdisplay{P(A \cap B) = P(A|B) \cdot P(B) = P(B|A) \cdot P(A)}

For three events:\PSdisplay{P(A \cap B \cap C) = P(A) \cdot P(B|A) \cdot P(C|A \cap B)}

\PSsection{1.8}{Independence}

\begin{PSbox}[unbreakable]{def}{Definition}

Events A and $B$ are \textbf{statistically independent} if:

\PSdisplay{P(A \cap B) = P(A) \cdot P(B)}

Equivalently: $P(A \mid B) = P(A)$ — knowing $B$ occurred gives no information about $A$.

\end{PSbox}

\PSsubsection{Engineering Significance}

Independence means \textbf{one event provides no information about the other}. This is a modeling assumption we often make:

\begin{itemize}
\tightlist
\item
  Noise samples at different times are independent
\item
  Bit errors in different time slots are independent (memoryless channel)
\item
  Component failures in different subsystems are independent
\end{itemize}

\PSsubsection{Pairwise vs. Mutual Independence}

For three events $A,\, B,\, C$:

\textbf{Pairwise independence} (all three pairs):

\begin{itemize}
\tightlist
\item
  $P(A \cap B) = P(A)P(B)$
\item
  $P(A \cap C) = P(A)P(C)$
\item
  $P(B \cap C) = P(B)P(C)$
\end{itemize}

\textbf{Mutual independence} (all combinations):

\begin{itemize}
\tightlist
\item
  All pairwise conditions above, AND
\item
  $P(A \cap B \cap C) = P(A)P(B)P(C)$
\end{itemize}

\begin{PSquote}

\PSwarn{} Pairwise independence does NOT imply mutual independence!

\end{PSquote}

\begin{PSbox}[unbreakable]{ex}{Engineering Example: Independent Component Failures}

A system has three independent modules, each with failure probability 0.1.

$P(\text{all three fail}) = P(A)P(B)P(C) = (0.1)^{3} = 0.001$

$P(\text{system works}) = P(\text{at least one works}) = 1 - P(\text{all fail}) = 1 - 0.001 = 0.999$

This calculation is ONLY valid because failures are independent.

\end{PSbox}

\PSsubsection{Testing Independence}

To verify independence from data:

\begin{enumerate}
\def\labelenumi{\arabic{enumi}.}
\tightlist
\item
  Compute $P(A),\, P(B)$, and $P(A \cap B)$ empirically
\item
  Check if $P(A \cap B) \approx P(A) \cdot P(B)$
\end{enumerate}

If not approximately equal, the events are \textbf{dependent}.

\PSsubsection{Common Misconception}

\begin{PSquote}

\PSno{} ``Mutually exclusive events are independent''

\end{PSquote}

This is \textbf{FALSE}. If $A \cap B = \emptyset$ and $P(A) > 0,\, P(B) > 0$, then:

\begin{itemize}
\tightlist
\item
  $P(A \cap B) = 0$
\item
  $P(A)P(B) > 0$
\item
  Therefore $P(A \cap B) \ne P(A)P(B)$
\end{itemize}

Mutually exclusive events are actually \textbf{maximally dependent} — if one occurs, the other cannot.

\PSsection{1.9}{Bayes’ Theorem}

\PSsubsection{Derivation}

From the definition of conditional probability:

\PSdisplay{P(A|B) = \frac{P(A \cap B)}{P(B)} = \frac{P(B|A) \cdot P(A)}{P(B)}}

\PSsubsection{Full Statement}

If $B_{1},\, B_{2},\, \ldots,\, B_{n}$ partition the sample space (mutually exclusive and exhaustive):

\PSdisplay{P(B_i|A) = \frac{P(A|B_i) \cdot P(B_i)}{\sum_{j=1}^{n} P(A|B_j) \cdot P(B_j)}}

\PSsubsection{Terminology}

{\def\LTcaptype{none} % do not increment counter
\begin{longtable}[c]{@{}
  >{\centering\arraybackslash}p{(0.965\linewidth - 4\tabcolsep) * \real{0.1818}}
  >{\centering\arraybackslash}p{(0.965\linewidth - 4\tabcolsep) * \real{0.1818}}
  >{\centering\arraybackslash}p{(0.965\linewidth - 4\tabcolsep) * \real{0.6364}}@{}}
\toprule\noalign{}
\rowcolor{PSnavy}\PSth{Term} & \PSth{Name} & \PSth{Engineering Meaning} \\
\midrule\noalign{}
\endhead
\bottomrule\noalign{}
\endlastfoot
$P(B_{i})$ & Prior probability & What we believed before observing data \\
$P(A \mid B_{i})$ & Likelihood & How likely is the observation given each hypothesis \\
$P(B_{i} \mid A)$ & Posterior probability & Updated belief after observing data \\
$P(A)$ & Evidence & Total probability of the observation \\
\end{longtable}
}

\begin{PSbox}[unbreakable]{ex}{Engineering Example: Fault Diagnosis}

A circuit board can have a fault in one of three subsystems:

\begin{itemize}
\tightlist
\item
  $B_{1}$ = power supply fault, $P(B_{1}) = 0.3$
\item
  $B_{2}$ = processor fault, $P(B_{2}) = 0.5$
\item
  $B_{3}$ = memory fault, $P(B_{3}) = 0.2$
\end{itemize}

A diagnostic test produces symptom $A$ (overheating):

\begin{itemize}
\tightlist
\item
  $P(A \mid B_{1}) = 0.8$ (power supply faults usually cause overheating)
\item
  $P(A \mid B_{2}) = 0.3$ (processor faults sometimes cause overheating)
\item
  $P(A \mid B_{3}) = 0.1$ (memory faults rarely cause overheating)
\end{itemize}

\textbf{Question:} Given that overheating is observed, what is the most likely fault?

\PSsolution{} 

$P(A) = P(A \mid B_{1})P(B_{1}) + P(A \mid B_{2})P(B_{2}) + P(A \mid B_{3})P(B_{3})$\PSbr{}$= (0.8)(0.3) + (0.3)(0.5) + (0.1)(0.2)$\PSbr{}$= 0.24 + 0.15 + 0.02 = 0.41$

$P(B_{1} \mid A) = \frac{(0.8)(0.3)}{0.41} = 0.585$ (most likely!)\PSbr{}$\displaystyle P(B_{2} \mid A) = \frac{(0.3)(0.5)}{0.41} = 0.366$\PSbr{}$\displaystyle P(B_{3} \mid A) = \frac{(0.1)(0.2)}{0.41} = 0.049$

\textbf{Engineering Decision:} Investigate the power supply first.

\end{PSbox}

\begin{PSbox}[unbreakable]{ex}{Engineering Example: Detection with Bayes}

Returning to the radar example from Section 1.7:

\begin{itemize}
\tightlist
\item
  $P(H_{1}) = 0.01,\, P(D \mid H_{1}) = 0.95,\, P(D \mid H_{0}) = 0.02$
\end{itemize}

$P(D) = P(D \mid H_{1})P(H_{1}) + P(D \mid H_{0})P(H_{0}) = (0.95)(0.01) + (0.02)(0.99) = 0.0095 + 0.0198 = 0.0293$

$\displaystyle P(H_{1} \mid D) = \frac{(0.95)(0.01)}{0.0293} = 0.324$

\textbf{Interpretation:} Even with a good detector (95\% detection, 2\% false alarm), when targets are rare (1\%), most detections are false alarms! Only 32.4\% of alarms correspond to real targets.

\textbf{Engineering Lesson:} System design must account for \textbf{prior probabilities}, not just detector performance.

\end{PSbox}

\PSsection{1.10}{Total Probability Theorem}

\PSsubsection{Statement}

If $B_{1},\, B_{2},\, \ldots,\, B_{n}$ form a \textbf{partition} of $S$ (mutually exclusive, exhaustive):

\PSdisplay{P(A) = \sum_{i=1}^{n} P(A|B_i) \cdot P(B_i)}

\PSsubsection{Interpretation}

To find the total probability of event A, decompose into all possible ``paths'' to A through the partition.

\begin{PSbox}[unbreakable]{ex}{Engineering Example: Multi-Path Communication}

A packet can be routed through three different paths:

\begin{itemize}
\tightlist
\item
  Path 1: chosen 50\% of the time, error probability 0.01
\item
  Path 2: chosen 30\% of the time, error probability 0.05
\item
  Path 3: chosen 20\% of the time, error probability 0.10
\end{itemize}

$P(\text{packet error}) = P(\text{error}\ \mid \text{path 1})P(\text{path 1}) + P(\text{error}\ \mid \text{path 2})P(\text{path 2}) + P(\text{error}\ \mid \text{path 3})P(\text{path 3})$\PSbr{}$= (0.01)(0.50) + (0.05)(0.30) + (0.10)(0.20)$\PSbr{}$= 0.005 + 0.015 + 0.020 = 0.040$

\textbf{Overall error rate: 4\%}

\end{PSbox}

\begin{PSbox}[unbreakable]{ex}{Engineering Example: Manufacturing with Multiple Suppliers}

A company sources capacitors from two suppliers:

\begin{itemize}
\tightlist
\item
  Supplier A: provides 60\% of capacitors, defect rate 2\%
\item
  Supplier $B$: provides 40\% of capacitors, defect rate 5\%
\end{itemize}

$P(\text{defective}) = P(\text{defective}\ \mid A)P(A) + P(\text{defective}\ \mid B)P(B)$\PSbr{}$= (0.02)(0.60) + (0.05)(0.40)$\PSbr{}$= 0.012 + 0.020 = 0.032$

If a defective capacitor is found, which supplier likely produced it?

$\displaystyle P(A \mid \text{defective}) = \frac{(0.02)(0.60)}{0.032} = 0.375$\PSbr{}$\displaystyle P(B \mid \text{defective}) = \frac{(0.05)(0.40)}{0.032} = 0.625$

Despite providing fewer units, Supplier $B$ is more likely responsible for a given defect.

\end{PSbox}

\PSsection{1.11}{MATLAB Examples}

\begin{PSbox}[unbreakable]{ex}{Example 1: Simulating Coin Flips (Random Experiment)}

\tcbset{PScodemode/.style={unbreakable}}

\begin{Shaded}
\begin{Highlighting}[]
\CommentTok{\%\% Simulating a Random Experiment: Biased Coin}
\CommentTok{\% A communication channel has bit error probability p = 0.1}
\CommentTok{\% Simulate transmitting N bits and count errors}

\VariableTok{p} \OperatorTok{=} \FloatTok{0.1}\OperatorTok{;}          \CommentTok{\% Bit error probability}
\VariableTok{N} \OperatorTok{=} \FloatTok{10000}\OperatorTok{;}        \CommentTok{\% Number of bits transmitted}

\CommentTok{\% Generate random outcomes: 1 = error, 0 = correct}
\VariableTok{errors} \OperatorTok{=} \VariableTok{rand}\NormalTok{(}\FloatTok{1}\OperatorTok{,} \VariableTok{N}\NormalTok{) }\OperatorTok{\textless{}} \VariableTok{p}\OperatorTok{;}

\CommentTok{\% Empirical probability of error}
\VariableTok{p\_empirical} \OperatorTok{=} \VariableTok{sum}\NormalTok{(}\VariableTok{errors}\NormalTok{) }\OperatorTok{/} \VariableTok{N}\OperatorTok{;}

\VariableTok{fprintf}\NormalTok{(}\SpecialStringTok{\textquotesingle{}Theoretical P(error) = \%.4f\textbackslash{}n\textquotesingle{}}\OperatorTok{,} \VariableTok{p}\NormalTok{)}\OperatorTok{;}
\VariableTok{fprintf}\NormalTok{(}\SpecialStringTok{\textquotesingle{}Empirical P(error)   = \%.4f (from \%d trials)\textbackslash{}n\textquotesingle{}}\OperatorTok{,} \VariableTok{p\_empirical}\OperatorTok{,} \VariableTok{N}\NormalTok{)}\OperatorTok{;}
\VariableTok{fprintf}\NormalTok{(}\SpecialStringTok{\textquotesingle{}Difference           = \%.4f\textbackslash{}n\textquotesingle{}}\OperatorTok{,} \VariableTok{abs}\NormalTok{(}\VariableTok{p} \OperatorTok{{-}} \VariableTok{p\_empirical}\NormalTok{))}\OperatorTok{;}
\end{Highlighting}
\end{Shaded}

\end{PSbox}

\begin{PSbox}[unbreakable]{ex}{Example 2: Verifying the Addition Rule}

\tcbset{PScodemode/.style={unbreakable}}

\begin{Shaded}
\begin{Highlighting}[]
\CommentTok{\%\% Verifying Inclusion{-}Exclusion with Simulation}
\CommentTok{\% Two independent failure modes:}
\CommentTok{\% A = overheating (P = 0.05)}
\CommentTok{\% B = power surge (P = 0.03)}

\VariableTok{N} \OperatorTok{=} \FloatTok{100000}\OperatorTok{;}       \CommentTok{\% Number of trials}
\VariableTok{pA} \OperatorTok{=} \FloatTok{0.05}\OperatorTok{;}
\VariableTok{pB} \OperatorTok{=} \FloatTok{0.03}\OperatorTok{;}

\CommentTok{\% Generate independent failures}
\VariableTok{A} \OperatorTok{=} \VariableTok{rand}\NormalTok{(}\FloatTok{1}\OperatorTok{,} \VariableTok{N}\NormalTok{) }\OperatorTok{\textless{}} \VariableTok{pA}\OperatorTok{;}
\VariableTok{B} \OperatorTok{=} \VariableTok{rand}\NormalTok{(}\FloatTok{1}\OperatorTok{,} \VariableTok{N}\NormalTok{) }\OperatorTok{\textless{}} \VariableTok{pB}\OperatorTok{;}

\CommentTok{\% Empirical probabilities}
\VariableTok{pA\_emp} \OperatorTok{=} \VariableTok{mean}\NormalTok{(}\VariableTok{A}\NormalTok{)}\OperatorTok{;}
\VariableTok{pB\_emp} \OperatorTok{=} \VariableTok{mean}\NormalTok{(}\VariableTok{B}\NormalTok{)}\OperatorTok{;}
\VariableTok{pAB\_emp} \OperatorTok{=} \VariableTok{mean}\NormalTok{(}\VariableTok{A} \OperatorTok{\&} \VariableTok{B}\NormalTok{)}\OperatorTok{;}           \CommentTok{\% Both occur}
\VariableTok{pAuB\_emp} \OperatorTok{=} \VariableTok{mean}\NormalTok{(}\VariableTok{A} \OperatorTok{|} \VariableTok{B}\NormalTok{)}\OperatorTok{;}          \CommentTok{\% At least one occurs}

\CommentTok{\% Theoretical (independent)}
\VariableTok{pAB\_theory} \OperatorTok{=} \VariableTok{pA} \OperatorTok{*} \VariableTok{pB}\OperatorTok{;}
\VariableTok{pAuB\_theory} \OperatorTok{=} \VariableTok{pA} \OperatorTok{+} \VariableTok{pB} \OperatorTok{{-}} \VariableTok{pA}\OperatorTok{*}\VariableTok{pB}\OperatorTok{;}

\VariableTok{fprintf}\NormalTok{(}\SpecialStringTok{\textquotesingle{}P(A ∩ B): Theory = \%.5f, Empirical = \%.5f\textbackslash{}n\textquotesingle{}}\OperatorTok{,} \VariableTok{pAB\_theory}\OperatorTok{,} \VariableTok{pAB\_emp}\NormalTok{)}\OperatorTok{;}
\VariableTok{fprintf}\NormalTok{(}\SpecialStringTok{\textquotesingle{}P(A ∪ B): Theory = \%.5f, Empirical = \%.5f\textbackslash{}n\textquotesingle{}}\OperatorTok{,} \VariableTok{pAuB\_theory}\OperatorTok{,} \VariableTok{pAuB\_emp}\NormalTok{)}\OperatorTok{;}
\end{Highlighting}
\end{Shaded}

\end{PSbox}

\begin{PSbox}[unbreakable]{ex}{Example 3: Conditional Probability Simulation}

\tcbset{PScodemode/.style={unbreakable}}

\begin{Shaded}
\begin{Highlighting}[]
\CommentTok{\%\% Conditional Probability: Binary Symmetric Channel}
\CommentTok{\% P(error | sent 0) = P(error | sent 1) = p = 0.1}

\VariableTok{p\_crossover} \OperatorTok{=} \FloatTok{0.1}\OperatorTok{;}
\VariableTok{N} \OperatorTok{=} \FloatTok{100000}\OperatorTok{;}

\CommentTok{\% Randomly generate transmitted bits (equally likely 0 or 1)}
\VariableTok{transmitted} \OperatorTok{=} \VariableTok{randi}\NormalTok{([}\FloatTok{0}\OperatorTok{,} \FloatTok{1}\NormalTok{]}\OperatorTok{,} \FloatTok{1}\OperatorTok{,} \VariableTok{N}\NormalTok{)}\OperatorTok{;}

\CommentTok{\% Channel: flip each bit with probability p\_crossover}
\VariableTok{noise} \OperatorTok{=} \VariableTok{rand}\NormalTok{(}\FloatTok{1}\OperatorTok{,} \VariableTok{N}\NormalTok{) }\OperatorTok{\textless{}} \VariableTok{p\_crossover}\OperatorTok{;}
\VariableTok{received} \OperatorTok{=} \VariableTok{xor}\NormalTok{(}\VariableTok{transmitted}\OperatorTok{,} \VariableTok{noise}\NormalTok{)}\OperatorTok{;}

\CommentTok{\% Compute conditional probabilities}
\VariableTok{sent\_0\_idx} \OperatorTok{=}\NormalTok{ (}\VariableTok{transmitted} \OperatorTok{==} \FloatTok{0}\NormalTok{)}\OperatorTok{;}
\VariableTok{sent\_1\_idx} \OperatorTok{=}\NormalTok{ (}\VariableTok{transmitted} \OperatorTok{==} \FloatTok{1}\NormalTok{)}\OperatorTok{;}

\CommentTok{\% P(received = 1 | sent = 0)}
\VariableTok{p\_1\_given\_sent0} \OperatorTok{=} \VariableTok{mean}\NormalTok{(}\VariableTok{received}\NormalTok{(}\VariableTok{sent\_0\_idx}\NormalTok{))}\OperatorTok{;}
\CommentTok{\% P(received = 0 | sent = 1)  }
\VariableTok{p\_0\_given\_sent1} \OperatorTok{=} \VariableTok{mean}\NormalTok{(}\VariableTok{received}\NormalTok{(}\VariableTok{sent\_1\_idx}\NormalTok{))}\OperatorTok{;}

\VariableTok{fprintf}\NormalTok{(}\SpecialStringTok{\textquotesingle{}P(receive 1 | sent 0) = \%.4f (theory: \%.4f)\textbackslash{}n\textquotesingle{}}\OperatorTok{,} \VariableTok{p\_1\_given\_sent0}\OperatorTok{,} \VariableTok{p\_crossover}\NormalTok{)}\OperatorTok{;}
\VariableTok{fprintf}\NormalTok{(}\SpecialStringTok{\textquotesingle{}P(receive 0 | sent 1) = \%.4f (theory: \%.4f)\textbackslash{}n\textquotesingle{}}\OperatorTok{,} \VariableTok{p\_0\_given\_sent1}\OperatorTok{,} \VariableTok{p\_crossover}\NormalTok{)}\OperatorTok{;}
\end{Highlighting}
\end{Shaded}

\end{PSbox}

\PSsubsection{Example 4: Bayes’ Theorem Simulation}

\tcbset{PScodemode/.style={unbreakable}}

\begin{Shaded}
\begin{Highlighting}[]
\CommentTok{\%\% Bayes\textquotesingle{} Theorem: Fault Diagnosis Simulation}
\VariableTok{N} \OperatorTok{=} \FloatTok{100000}\OperatorTok{;}

\CommentTok{\% Prior probabilities of fault types}
\VariableTok{p\_power} \OperatorTok{=} \FloatTok{0.3}\OperatorTok{;}   \VariableTok{p\_proc} \OperatorTok{=} \FloatTok{0.5}\OperatorTok{;}   \VariableTok{p\_mem} \OperatorTok{=} \FloatTok{0.2}\OperatorTok{;}

\CommentTok{\% Likelihoods of symptom (overheating) given fault type}
\VariableTok{p\_sym\_power} \OperatorTok{=} \FloatTok{0.8}\OperatorTok{;}  \VariableTok{p\_sym\_proc} \OperatorTok{=} \FloatTok{0.3}\OperatorTok{;}  \VariableTok{p\_sym\_mem} \OperatorTok{=} \FloatTok{0.1}\OperatorTok{;}

\CommentTok{\% Simulate: generate fault type for each trial}
\VariableTok{fault\_type} \OperatorTok{=} \VariableTok{zeros}\NormalTok{(}\FloatTok{1}\OperatorTok{,} \VariableTok{N}\NormalTok{)}\OperatorTok{;}
\VariableTok{r} \OperatorTok{=} \VariableTok{rand}\NormalTok{(}\FloatTok{1}\OperatorTok{,} \VariableTok{N}\NormalTok{)}\OperatorTok{;}
\VariableTok{fault\_type}\NormalTok{(}\VariableTok{r} \OperatorTok{\textless{}} \VariableTok{p\_power}\NormalTok{) }\OperatorTok{=} \FloatTok{1}\OperatorTok{;}                          \CommentTok{\% Power supply}
\VariableTok{fault\_type}\NormalTok{(}\VariableTok{r} \OperatorTok{\textgreater{}=} \VariableTok{p\_power} \OperatorTok{\&} \VariableTok{r} \OperatorTok{\textless{}} \VariableTok{p\_power} \OperatorTok{+} \VariableTok{p\_proc}\NormalTok{) }\OperatorTok{=} \FloatTok{2}\OperatorTok{;}  \CommentTok{\% Processor}
\VariableTok{fault\_type}\NormalTok{(}\VariableTok{r} \OperatorTok{\textgreater{}=} \VariableTok{p\_power} \OperatorTok{+} \VariableTok{p\_proc}\NormalTok{) }\OperatorTok{=} \FloatTok{3}\OperatorTok{;}                \CommentTok{\% Memory}

\CommentTok{\% Generate symptom based on fault type}
\VariableTok{symptom} \OperatorTok{=} \VariableTok{zeros}\NormalTok{(}\FloatTok{1}\OperatorTok{,} \VariableTok{N}\NormalTok{)}\OperatorTok{;}
\VariableTok{symptom}\NormalTok{(}\VariableTok{fault\_type} \OperatorTok{==} \FloatTok{1}\NormalTok{) }\OperatorTok{=} \VariableTok{rand}\NormalTok{(}\FloatTok{1}\OperatorTok{,} \VariableTok{sum}\NormalTok{(}\VariableTok{fault\_type}\OperatorTok{==}\FloatTok{1}\NormalTok{)) }\OperatorTok{\textless{}} \VariableTok{p\_sym\_power}\OperatorTok{;}
\VariableTok{symptom}\NormalTok{(}\VariableTok{fault\_type} \OperatorTok{==} \FloatTok{2}\NormalTok{) }\OperatorTok{=} \VariableTok{rand}\NormalTok{(}\FloatTok{1}\OperatorTok{,} \VariableTok{sum}\NormalTok{(}\VariableTok{fault\_type}\OperatorTok{==}\FloatTok{2}\NormalTok{)) }\OperatorTok{\textless{}} \VariableTok{p\_sym\_proc}\OperatorTok{;}
\VariableTok{symptom}\NormalTok{(}\VariableTok{fault\_type} \OperatorTok{==} \FloatTok{3}\NormalTok{) }\OperatorTok{=} \VariableTok{rand}\NormalTok{(}\FloatTok{1}\OperatorTok{,} \VariableTok{sum}\NormalTok{(}\VariableTok{fault\_type}\OperatorTok{==}\FloatTok{3}\NormalTok{)) }\OperatorTok{\textless{}} \VariableTok{p\_sym\_mem}\OperatorTok{;}

\CommentTok{\% Compute posterior: P(fault\_type | symptom = 1)}
\VariableTok{sym\_idx} \OperatorTok{=} \VariableTok{logical}\NormalTok{(}\VariableTok{symptom}\NormalTok{)}\OperatorTok{;}
\VariableTok{p\_power\_given\_sym} \OperatorTok{=} \VariableTok{mean}\NormalTok{(}\VariableTok{fault\_type}\NormalTok{(}\VariableTok{sym\_idx}\NormalTok{) }\OperatorTok{==} \FloatTok{1}\NormalTok{)}\OperatorTok{;}
\VariableTok{p\_proc\_given\_sym} \OperatorTok{=} \VariableTok{mean}\NormalTok{(}\VariableTok{fault\_type}\NormalTok{(}\VariableTok{sym\_idx}\NormalTok{) }\OperatorTok{==} \FloatTok{2}\NormalTok{)}\OperatorTok{;}
\VariableTok{p\_mem\_given\_sym} \OperatorTok{=} \VariableTok{mean}\NormalTok{(}\VariableTok{fault\_type}\NormalTok{(}\VariableTok{sym\_idx}\NormalTok{) }\OperatorTok{==} \FloatTok{3}\NormalTok{)}\OperatorTok{;}

\CommentTok{\% Theoretical values}
\VariableTok{p\_sym} \OperatorTok{=} \VariableTok{p\_sym\_power}\OperatorTok{*}\VariableTok{p\_power} \OperatorTok{+} \VariableTok{p\_sym\_proc}\OperatorTok{*}\VariableTok{p\_proc} \OperatorTok{+} \VariableTok{p\_sym\_mem}\OperatorTok{*}\VariableTok{p\_mem}\OperatorTok{;}
\VariableTok{theory\_power} \OperatorTok{=} \VariableTok{p\_sym\_power}\OperatorTok{*}\VariableTok{p\_power} \OperatorTok{/} \VariableTok{p\_sym}\OperatorTok{;}
\VariableTok{theory\_proc} \OperatorTok{=} \VariableTok{p\_sym\_proc}\OperatorTok{*}\VariableTok{p\_proc} \OperatorTok{/} \VariableTok{p\_sym}\OperatorTok{;}
\VariableTok{theory\_mem} \OperatorTok{=} \VariableTok{p\_sym\_mem}\OperatorTok{*}\VariableTok{p\_mem} \OperatorTok{/} \VariableTok{p\_sym}\OperatorTok{;}

\VariableTok{fprintf}\NormalTok{(}\SpecialStringTok{\textquotesingle{}P(power fault | symptom): Empirical = \%.3f, Theory = \%.3f\textbackslash{}n\textquotesingle{}}\OperatorTok{,} \VariableTok{p\_power\_given\_sym}\OperatorTok{,} \VariableTok{theory\_power}\NormalTok{)}\OperatorTok{;}
\VariableTok{fprintf}\NormalTok{(}\SpecialStringTok{\textquotesingle{}P(proc fault  | symptom): Empirical = \%.3f, Theory = \%.3f\textbackslash{}n\textquotesingle{}}\OperatorTok{,} \VariableTok{p\_proc\_given\_sym}\OperatorTok{,} \VariableTok{theory\_proc}\NormalTok{)}\OperatorTok{;}
\VariableTok{fprintf}\NormalTok{(}\SpecialStringTok{\textquotesingle{}P(mem fault   | symptom): Empirical = \%.3f, Theory = \%.3f\textbackslash{}n\textquotesingle{}}\OperatorTok{,} \VariableTok{p\_mem\_given\_sym}\OperatorTok{,} \VariableTok{theory\_mem}\NormalTok{)}\OperatorTok{;}
\end{Highlighting}
\end{Shaded}

\begin{PSbox}[unbreakable]{ex}{Example 5: Independence Test}

\tcbset{PScodemode/.style={unbreakable}}

\begin{Shaded}
\begin{Highlighting}[]
\CommentTok{\%\% Testing Independence vs. Dependence}
\VariableTok{N} \OperatorTok{=} \FloatTok{100000}\OperatorTok{;}

\CommentTok{\% Case 1: Independent events}
\VariableTok{A\_ind} \OperatorTok{=} \VariableTok{rand}\NormalTok{(}\FloatTok{1}\OperatorTok{,} \VariableTok{N}\NormalTok{) }\OperatorTok{\textless{}} \FloatTok{0.3}\OperatorTok{;}
\VariableTok{B\_ind} \OperatorTok{=} \VariableTok{rand}\NormalTok{(}\FloatTok{1}\OperatorTok{,} \VariableTok{N}\NormalTok{) }\OperatorTok{\textless{}} \FloatTok{0.4}\OperatorTok{;}  \CommentTok{\% Generated independently}

\VariableTok{pA} \OperatorTok{=} \VariableTok{mean}\NormalTok{(}\VariableTok{A\_ind}\NormalTok{)}\OperatorTok{;}
\VariableTok{pB} \OperatorTok{=} \VariableTok{mean}\NormalTok{(}\VariableTok{B\_ind}\NormalTok{)}\OperatorTok{;}
\VariableTok{pAB} \OperatorTok{=} \VariableTok{mean}\NormalTok{(}\VariableTok{A\_ind} \OperatorTok{\&} \VariableTok{B\_ind}\NormalTok{)}\OperatorTok{;}
\VariableTok{fprintf}\NormalTok{(}\SpecialStringTok{\textquotesingle{}Independent case:\textbackslash{}n\textquotesingle{}}\NormalTok{)}\OperatorTok{;}
\VariableTok{fprintf}\NormalTok{(}\SpecialStringTok{\textquotesingle{}  P(A)*P(B) = \%.4f, P(A∩B) = \%.4f\textbackslash{}n\textquotesingle{}}\OperatorTok{,} \VariableTok{pA}\OperatorTok{*}\VariableTok{pB}\OperatorTok{,} \VariableTok{pAB}\NormalTok{)}\OperatorTok{;}

\CommentTok{\% Case 2: Dependent events (B depends on A)}
\VariableTok{A\_dep} \OperatorTok{=} \VariableTok{rand}\NormalTok{(}\FloatTok{1}\OperatorTok{,} \VariableTok{N}\NormalTok{) }\OperatorTok{\textless{}} \FloatTok{0.3}\OperatorTok{;}
\VariableTok{B\_dep} \OperatorTok{=} \VariableTok{zeros}\NormalTok{(}\FloatTok{1}\OperatorTok{,} \VariableTok{N}\NormalTok{)}\OperatorTok{;}
\VariableTok{B\_dep}\NormalTok{(}\VariableTok{A\_dep} \OperatorTok{==} \FloatTok{1}\NormalTok{) }\OperatorTok{=} \VariableTok{rand}\NormalTok{(}\FloatTok{1}\OperatorTok{,} \VariableTok{sum}\NormalTok{(}\VariableTok{A\_dep}\NormalTok{)) }\OperatorTok{\textless{}} \FloatTok{0.7}\OperatorTok{;}   \CommentTok{\% P(B|A) = 0.7}
\VariableTok{B\_dep}\NormalTok{(}\VariableTok{A\_dep} \OperatorTok{==} \FloatTok{0}\NormalTok{) }\OperatorTok{=} \VariableTok{rand}\NormalTok{(}\FloatTok{1}\OperatorTok{,} \VariableTok{sum}\NormalTok{(}\OperatorTok{\textasciitilde{}}\VariableTok{A\_dep}\NormalTok{)) }\OperatorTok{\textless{}} \FloatTok{0.2}\OperatorTok{;}  \CommentTok{\% P(B|not A) = 0.2}

\VariableTok{pA} \OperatorTok{=} \VariableTok{mean}\NormalTok{(}\VariableTok{A\_dep}\NormalTok{)}\OperatorTok{;}
\VariableTok{pB} \OperatorTok{=} \VariableTok{mean}\NormalTok{(}\VariableTok{B\_dep}\NormalTok{)}\OperatorTok{;}
\VariableTok{pAB} \OperatorTok{=} \VariableTok{mean}\NormalTok{(}\VariableTok{A\_dep} \OperatorTok{\&} \VariableTok{B\_dep}\NormalTok{)}\OperatorTok{;}
\VariableTok{fprintf}\NormalTok{(}\SpecialStringTok{\textquotesingle{}Dependent case:\textbackslash{}n\textquotesingle{}}\NormalTok{)}\OperatorTok{;}
\VariableTok{fprintf}\NormalTok{(}\SpecialStringTok{\textquotesingle{}  P(A)*P(B) = \%.4f, P(A∩B) = \%.4f\textbackslash{}n\textquotesingle{}}\OperatorTok{,} \VariableTok{pA}\OperatorTok{*}\VariableTok{pB}\OperatorTok{,} \VariableTok{pAB}\NormalTok{)}\OperatorTok{;}
\VariableTok{fprintf}\NormalTok{(}\SpecialStringTok{\textquotesingle{}  These differ =\textgreater{} events are dependent\textbackslash{}n\textquotesingle{}}\NormalTok{)}\OperatorTok{;}
\end{Highlighting}
\end{Shaded}

\end{PSbox}

\PSsection{1.12}{Summary: Key Distinctions}

{\def\LTcaptype{none} % do not increment counter
\begin{longtable}[c]{@{}
  >{\centering\arraybackslash}p{(0.965\linewidth - 4\tabcolsep) * \real{0.2195}}
  >{\centering\arraybackslash}p{(0.965\linewidth - 4\tabcolsep) * \real{0.2683}}
  >{\centering\arraybackslash}p{(0.965\linewidth - 4\tabcolsep) * \real{0.5122}}@{}}
\toprule\noalign{}
\rowcolor{PSnavy}\PSth{Concept} & \PSth{Definition} & \PSth{Engineering Meaning} \\
\midrule\noalign{}
\endhead
\bottomrule\noalign{}
\endlastfoot
\textbf{Event} & A subset of the sample space & A specific outcome or set of outcomes we care about \\
\textbf{Probability} & $P(A)$ = measure of likelihood & Long-run relative frequency of occurrence \\
\textbf{Conditional Probability} & $P(A \mid B)$ = updated probability given $B$ & How knowledge of one event changes another’s likelihood \\
\textbf{Independence} & $P(A \cap B) = P(A)P(B)$ & Knowing one event tells nothing about the other \\
\end{longtable}
}

\PSsection{1.13}{Practice Problems}

\begin{PSbox}[unbreakable]{prob}{Problem 1: Component Reliability}

A system has two independent processors. Processor 1 fails with probability 0.02, Processor 2 fails with probability 0.03. The system fails only if BOTH processors fail.

\begin{enumerate}
\def\labelenumi{(\alph{enumi})}
\tightlist
\item
  What is the probability that the system fails?
\item
  What is the probability that at least one processor fails?
\item
  Given that the system is still working, what is the probability that processor 1 has failed?
\end{enumerate}

\PSsolution{} 

\begin{enumerate}
\def\labelenumi{(\alph{enumi})}
\item
  $P(\text{system fails}) = P(\text{both fail}) = (0.02)(0.03) = 0.0006$
\item
  $P(\text{at least one fails}) = 1 - P(\text{none fail}) = 1 - (0.98)(0.97) = 1 - 0.9506 = 0.0494$
\item
  System working = NOT(both fail) = at least one works.\PSbr{}$\displaystyle P(P1\ \text{failed}\ \mid \text{system works}) = \frac{P(P1\ \text{failed}\ \cap P2\ \text{works})}{P(\text{system works})}$\PSbr{}$\displaystyle = \frac{(0.02)(0.97)}{1 - 0.0006} = \frac{0.0194}{0.9994} \approx 0.0194$
\end{enumerate}

\end{PSbox}

\begin{PSbox}[unbreakable]{prob}{Problem 2: Sensor Detection}

A fire detection system has: $P(\text{fire}) = 0.001,\, P(\text{alarm}\ \mid \text{fire}) = 0.98,\, P(\text{alarm}\ \mid \text{no fire}) = 0.01$.

\begin{enumerate}
\def\labelenumi{(\alph{enumi})}
\tightlist
\item
  What is $P(\text{alarm})$?
\item
  If the alarm sounds, what is $P(\text{fire}\ \mid \text{alarm})$?
\item
  Comment on the result.
\end{enumerate}

\PSsolution{} 

\begin{enumerate}
\def\labelenumi{(\alph{enumi})}
\item
  $P(\text{alarm}) = P(\text{alarm}\ \mid \text{fire})P(\text{fire}) + P(\text{alarm}\ \mid \text{no fire})P(\text{no fire})$\PSbr{}$= (0.98)(0.001) + (0.01)(0.999) = 0.000980 + 0.009990 = 0.010970$
\item
  $\displaystyle P(\text{fire}\ \mid \text{alarm}) = \frac{P(\text{alarm}\ \mid \text{fire})P(\text{fire})}{P(\text{alarm})}$\PSbr{}$\displaystyle = \frac{(0.98)(0.001)}{0.01097} = 0.0894$
\item
  Only $\sim 9\%$ of alarms correspond to real fires! This is the base rate fallacy. Even with an excellent detector, rare events produce many false alarms.
\end{enumerate}

\end{PSbox}

\begin{PSbox}[unbreakable]{prob}{Problem 3: Communication Error}

A packet of 8 bits is transmitted. Each bit has independent error probability 0.05.

\begin{enumerate}
\def\labelenumi{(\alph{enumi})}
\tightlist
\item
  What is $P(\text{no errors in the packet})$?
\item
  What is $P(\text{exactly one error})$?
\item
  What is $P(\text{at least one error})$?
\item
  Write MATLAB code to verify by simulation.
\end{enumerate}

\PSsolution{} 

\begin{enumerate}
\def\labelenumi{(\alph{enumi})}
\item
  $P(\text{no errors}) = (0.95)^{8} = 0.6634$
\item
  $P(\text{exactly one error}) = C(8,\,1)(0.05)^{1}(0.95)^{7} = 8 \times 0.05 \times 0.6983 = 0.2793$
\item
  $P(\text{at least one error}) = 1 - P(\text{no errors}) = 1 - 0.6634 = 0.3366$
\item
  MATLAB verification:
\end{enumerate}

\tcbset{PScodemode/.style={unbreakable}}

\begin{Shaded}
\begin{Highlighting}[]
\VariableTok{p} \OperatorTok{=} \FloatTok{0.05}\OperatorTok{;} \VariableTok{n\_bits} \OperatorTok{=} \FloatTok{8}\OperatorTok{;} \VariableTok{N\_packets} \OperatorTok{=} \FloatTok{100000}\OperatorTok{;}
\VariableTok{errors\_per\_packet} \OperatorTok{=} \VariableTok{sum}\NormalTok{(}\VariableTok{rand}\NormalTok{(}\VariableTok{n\_bits}\OperatorTok{,} \VariableTok{N\_packets}\NormalTok{) }\OperatorTok{\textless{}} \VariableTok{p}\OperatorTok{,} \FloatTok{1}\NormalTok{)}\OperatorTok{;}
\VariableTok{fprintf}\NormalTok{(}\SpecialStringTok{\textquotesingle{}P(0 errors) = \%.4f (theory: \%.4f)\textbackslash{}n\textquotesingle{}}\OperatorTok{,} \VariableTok{mean}\NormalTok{(}\VariableTok{errors\_per\_packet}\OperatorTok{==}\FloatTok{0}\NormalTok{)}\OperatorTok{,} \FloatTok{0.95}\OperatorTok{\^{}}\FloatTok{8}\NormalTok{)}\OperatorTok{;}
\VariableTok{fprintf}\NormalTok{(}\SpecialStringTok{\textquotesingle{}P(1 error)  = \%.4f (theory: \%.4f)\textbackslash{}n\textquotesingle{}}\OperatorTok{,} \VariableTok{mean}\NormalTok{(}\VariableTok{errors\_per\_packet}\OperatorTok{==}\FloatTok{1}\NormalTok{)}\OperatorTok{,} \FloatTok{8}\OperatorTok{*}\FloatTok{0.05}\OperatorTok{*}\FloatTok{0.95}\OperatorTok{\^{}}\FloatTok{7}\NormalTok{)}\OperatorTok{;}
\VariableTok{fprintf}\NormalTok{(}\SpecialStringTok{\textquotesingle{}P(\textgreater{}=1 error)= \%.4f (theory: \%.4f)\textbackslash{}n\textquotesingle{}}\OperatorTok{,} \VariableTok{mean}\NormalTok{(}\VariableTok{errors\_per\_packet}\OperatorTok{\textgreater{}=}\FloatTok{1}\NormalTok{)}\OperatorTok{,} \FloatTok{1}\OperatorTok{{-}}\FloatTok{0.95}\OperatorTok{\^{}}\FloatTok{8}\NormalTok{)}\OperatorTok{;}
\end{Highlighting}
\end{Shaded}

\end{PSbox}

\begin{PSbox}[unbreakable]{prob}{Problem 4: Total Probability Application}

A wireless system operates in three modes:

\begin{itemize}
\tightlist
\item
  Mode 1 (good channel): 60\% of time, $P(\text{outage}) = 0.01$
\item
  Mode 2 (moderate channel): 30\% of time, $P(\text{outage}) = 0.10$
\item
  Mode 3 (poor channel): 10\% of time, $P(\text{outage}) = 0.50$
\end{itemize}

\begin{enumerate}
\def\labelenumi{(\alph{enumi})}
\tightlist
\item
  What is the overall outage probability?
\item
  Given an outage occurred, what is the probability the system was in Mode 3?
\end{enumerate}

\PSsolution{} 

\begin{enumerate}
\def\labelenumi{(\alph{enumi})}
\item
  $P(\text{outage}) = (0.01)(0.60) + (0.10)(0.30) + (0.50)(0.10) = 0.006 + 0.030 + 0.050 = 0.086$
\item
  $\displaystyle P(\text{Mode 3}\ \mid \text{outage}) = \frac{(0.50)(0.10)}{0.086} = 0.581$
\end{enumerate}

\end{PSbox}

\begin{PSbox}[unbreakable]{prob}{Problem 5: Independence Analysis}

Two sensors monitor a process. Sensor A has false alarm probability 0.05. Sensor $B$ has false alarm probability 0.03.

\begin{enumerate}
\def\labelenumi{(\alph{enumi})}
\tightlist
\item
  If sensors are independent, what is $P(\text{both give false alarm})$?
\item
  If sensors share a noise source and $P(\text{both false alarm}) = 0.01$, are they independent?
\item
  What is $P(A\ \text{false alarm}\ \mid B\ \text{false alarm})$ in case (b)?
\end{enumerate}

\PSsolution{} 

\begin{enumerate}
\def\labelenumi{(\alph{enumi})}
\item
  If independent: $P(\text{both}) = (0.05)(0.03) = 0.0015$
\item
  If $P(\text{both}) = 0.01 \ne 0.0015$, they are NOT independent. The shared noise source creates dependence.
\item
  $\displaystyle P(A \mid B) = \frac{P(A \cap B)}{P(B)} = \frac{0.01}{0.03} = 0.333$
\end{enumerate}

This is much higher than $P(A) = 0.05$, confirming strong dependence.

\end{PSbox}

\PSsection{1.14}{Key Takeaways}

\begin{enumerate}
\def\labelenumi{\arabic{enumi}.}
\tightlist
\item
  \textbf{Probability quantifies uncertainty} — essential for engineering design under uncertainty
\item
  \textbf{Conditional probability updates beliefs} — the foundation of detection, estimation, and diagnosis
\item
  \textbf{Independence simplifies analysis} — but must be justified, not assumed
\item
  \textbf{Bayes’ theorem inverts conditional probabilities} — from ``likelihood of evidence given cause'' to ``likelihood of cause given evidence''
\item
  \textbf{Base rates matter} — even good detectors produce false alarms when events are rare
\item
  \textbf{Simulation validates analysis} — MATLAB Monte Carlo confirms theoretical calculations
\end{enumerate}

\emph{Next Module: {Module 2 — One Random Variable and Its Functions}}
